\section{Conclusions}
\label{sec:conclusion}
RAG models, with their compact size and personalized responses, are well-suited for edge deployment. To address system-level challenges in this context, we dissect RAG pipeline to quantify how various factors impact performance on edge devices. Through detailed stage-wise characterization, we identify inefficiencies in existing solutions and leverage these insights to design a multi-width, resource-aware pipeline that better utilizes available compute. Our implementation, MaestroRAG, targets two common edge use cases: latency-critical and throughput-critical tasks. Without modifying hardware or model architecture, our hardware-aware, software-only scheme achieves up to $12\times$ latency reduction and $5.6\times$ throughput improvement over state-of-the-art baselines. This makes RAG more viable for edge deployment.

\section{AI Usage Disclosure}
Generative AI tools were used for grammar check and sentence formatting during manuscript preparation. Agentic coding tools were used for graph preparation and experimental automation. All technical content, experimental results, algorithms, and architectural designs are the original work of the authors.
% RAG models, due to their compact size and personalized response, are the primary candidates for edge deployment.
% To address the system-level intricacies of RAG on edge devices, in this paper, we dissect into the RAG pipeline to gain insights and identify the factors and the degree to which they affect the performance of RAG applications on edge devices. With detailed characterization of each individual stage, we identify the suboptimality of the existing solutions. Furthermore, we extend our insights to propose a multi-width resource-aware RAG pipeline which utilizes the compute resources available to it better than the existing solutions and achieves speedups. We develop our proposed solution, \design{}, to accommodate two most frequently occurring use cases for edge users: latency-critical tasks and throughput-critical tasks. Without requiring any changes in the underlying hardware or model, our proposed hardware-aware, software-only scheme achieves up to $12\times$ improvement in latency and up to $6.7\times$ improvement in throughput over the state-of-the-art. This makes the deployment of RAG more suitable in edge devices and creates a new baseline for future research work.
%